\documentclass[12pt,aspectratio=43]{beamer}

% 自訂字體的封包
\usepackage{fontspec}

% 導入圖形與表格的封包
\usepackage{graphicx}
\graphicspath{{./20251203_Meeting}}

\usepackage{booktabs}

% 排列多個子圖形的封包
\usepackage{subfigure} 

 % 提供進階數學環境與公式排版
\usepackage{mathtools, amsmath, amsfonts, amsthm, latexsym} 

%% 設定中文字體
\usepackage{xeCJK}
\setCJKmainfont[AutoFakeBold=3.5 , AutoFakeSlant=0.2]{標楷體}
\setmainfont{Times New Roman}
\XeTeXlinebreaklocale "zh"
\XeTeXlinebreakskip = 0pt plus 1pt

% 自訂字體顏色的封包
\usepackage{color} 

% 圖表自動編號的封包
\usepackage{caption} 

% 允許表格的一格能多列呈現的封包
\usepackage{multirow} 

% 可指定表格排版的封包
\usepackage{array}

% 翻轉表格的封包
\usepackage{lscape} 

% 序列標號
\usepackage{enumerate} 

%% 設定自動編號
\setbeamertemplate{caption}[numbered]

%% 自訂顏色
\definecolor{Myblue}{RGB}{57,108,144}
\definecolor{Default_Blue}{RGB}{52,51,171}

% (動畫)字會若隱若現的顯示
\beamersetuncovermixins{\opaqueness<1->{45}}{\opaqueness<1->{95}} 

% 設定中文的標籤
\renewcommand{\figurename}{圖} 
\renewcommand{\tablename}{表} 

% 排版形式 
\mode<presentation>{
\usetheme{Berlin}
% enumerate 及 item 的形狀
\setbeamertemplate{itemize items}[triangle]
% 調整外框形式的字體大小
\setbeamerfont{headline}{size=\scriptsize}
\setbeamerfont{footline}{size=\scriptsize}
% 取消右下方的跳轉工具列
\setbeamertemplate{navigation symbols}{} 
% 自定義2：清除外框下緣但僅出頁碼
\setbeamertemplate{footline}[page number] 
% 顏色主題
\usecolortheme{dove}
% 設定頁面邊界
\setbeamersize{text margin left=0.6cm, text margin right=0.6cm}
\special{papersize=\the\paperwidth,\the\paperheight}
\providecommand{\tabularnewline}{\\}
}


\title{\fontsize{24pt}{20pt}統計諮詢與專業實務\ 期末}
\author{\fontsize{12pt}{0pt}611211106 黃丞暘\\ 611311005 吳沁璇\\ 611411104 蔡旻宏}
\institute{\fontsize{12pt}{0pt}國立東華大學應用數學所}
\date{December 11, 2025}

\begin{document}

% 封面
\begin{frame}
  \titlepage
\end{frame}

% 目錄
\begin{frame}{\textbf{目錄}}
  \tableofcontents
\end{frame}

% 研究動機
\section{研究動機}

\begin{frame}{\textbf{研究動機}}
\linespread{1.5}
  {輸入文字}
  \begin{itemize}
    \item
  \end{itemize}
  \begin{itemize}
    \item 
  \end{itemize}
\end{frame}

% 研究方法
\section{研究方法}

\begin{frame}{\textbf{研究方法}}
\linespread{1.5}
  {輸入文字}
  \begin{itemize}
    \item
  \end{itemize}
  \begin{itemize}
    \item 
  \end{itemize}
\end{frame}

% 研究流程
\section{研究流程}

\begin{frame}{\textbf{研究流程}}
\linespread{1.5}
  {輸入文字}
  \begin{itemize}
    \item
  \end{itemize}
  \begin{itemize}
    \item 
  \end{itemize}
\end{frame}

% 研究資料
\section{研究資料}

\begin{frame}{\textbf{研究資料}}
\linespread{1.5}
  {輸入文字}
  \begin{itemize}
    \item
  \end{itemize}
  \begin{itemize}
    \item 
  \end{itemize}
\end{frame}

% 研究結果
\section{研究結果}

\begin{frame}{\textbf{研究結果}}
\linespread{1.5}
  {輸入文字}
  \begin{itemize}
    \item
  \end{itemize}
  \begin{itemize}
    \item 
  \end{itemize}
\end{frame}

% 結論
\section{結論}

\begin{frame}{\textbf{結論}}
\linespread{1.5}
  {輸入文字}
  \begin{itemize}
    \item
  \end{itemize}
  \begin{itemize}
    \item 
  \end{itemize}
\end{frame}

% END
\begin{frame}[plain] % 產生無外框的頁面
	\begin{center}
		{\Huge 結束　謝謝聆聽}
		\bigskip\bigskip 
		
		{\LARGE 問題}
	\end{center}
\end{frame}
\end{document}
